\section{Conclusion}

% We introduced the RLPL framework, a novel approach that maximizes the utility of unlabeled data. RLPL addresses key limitations in current methods, such as dependence on large batch sizes and minimal unlabeled data utilization in downstream tasks. Our experiments show that RLPL consistently outperforms other SOTA methods, achieving superior accuracy with reduced batch sizes, making it an efficient choice for memory-constrained environments. RLPL achieved significant improvements across various datasets, demonstrating a notable average performance gain, especially in scenarios with limited labeled data. Beyond accuracy gains, RLPL exhibits strong transferability, demonstrated by its high performance across a range of diverse downstream classification tasks without requiring pseudo labels for adaptation. Overall, RLPL offers a significant improvement in semi-supervised learning, with enhanced data efficiency, robustness, and transferability, making it ideal for diverse memory-efficient applications and strong cross-task generalization.

We present RLPL, a novel semi-supervised learning framework designed to maximize the utility of unlabeled data, addressing critical limitations in existing methods such as the reliance on large batch sizes and underutilization of unlabeled data in downstream tasks. RLPL demonstrates consistent improvements over SOTA methods, achieving high accuracy with smaller batch sizes, which enhances efficiency for memory-constrained environments. Our experiments show that RLPL yields substantial performance gains across multiple datasets, with notable advantages in settings with limited labeled data. In addition to accuracy improvements, RLPL demonstrates exceptional transferability, achieving strong performance across a diverse range of downstream classification tasks without needing pseudo labels for adaptation. By offering enhanced data efficiency, robustness, and generalization, RLPL is ideally suited for a variety of applications that demand memory efficiency and cross-task adaptability, marking a significant advancement in semi-supervised learning.